\subsection{Finite-order estimates for the mean}\label{cc:uf:mean:sec:finite}
The phase-independent corrections of the expansion solve the linearized
Euler equations over the parent, with zero displacement at both ends of
$[0,T]$. For a smooth phase-independent source \(\bar f(t,a)\),
Theorem~\ref{cc:uf:mean:inverse} over $U$
gives \((B,\bar q)\) with
\begin{equation}\label{cc:uf:hier:eq:mean}
 B_t+MB+\mathbf d\bar q=\bar f,\qquad \mathbf d\cdot B=0,\qquad
 z(0)=z(T)=0,\quad B=D_aX z_t.
\end{equation}
In the notation of that theorem, $\bar f=f\circ X$,
$g=(D_aX)^T\bar f$, $v=z_t$, and $\bar q$ is the pressure $\mathsf q$ of
\eqref{cc:uf:mean:strong}. We call \eqref{cc:uf:hier:eq:mean} the
\emph{mean equation}. The terminal condition selects the initial
velocity $B(0)$, so the mean can already be present before activation.
We estimate it on all of $[0,T]$ in the unweighted label norms.
The inverse preserves the derivative order in the scaled H\"older
norms. Passing to the supremum norms by \eqref{cc:uf:trunc:eq:holder-sup}
costs one derivative of the source. At each level of the expansion
the mean source is formed from terms already estimated at all orders.

\begin{proposition}\label{cc:uf:mean:columns}
There is an absolute exponent $C$ such that, for every $n\ge0$, every
solution of the mean equation \eqref{cc:uf:hier:eq:mean} satisfies
\begin{equation}\label{cc:uf:mean:bound}
 \norm B_{[n]}\le C_n\mathsf D^C\norm{\bar f}_{[n+1]}.
\end{equation}
\end{proposition}
\begin{proof}
By Proposition~\ref{cc:uf:trunc:prop:parent} and
Lemma~\ref{cc:uf:trunc:lem:coefficients}(a), the coefficients $K$ and
$\Omega$ of \eqref{cc:uf:mean:data}, which enter
\eqref{cc:uf:mean:integrated}, satisfy \eqref{cc:uf:trunc:eq:power}.
Theorem~\ref{cc:uf:mean:inverse} applies over $U$ with $D_-=\mathsf D$,
by \eqref{cc:uf:tower:eq:parent-bound}. Its constant $A$ of
\eqref{cc:uf:mean:base-constant} contains $D_-$,
$\|M\|_\infty\le\mathsf D$ and the norms of $K$ and $\Omega$ in
$C_t\Holder{1}$, so it is at most $\mathsf D^C$. Since $\ell_*$ is a
sufficiently large negative power of $\mathsf D$, the first bound of
\eqref{cc:uf:trunc:eq:holder-sup} and
Lemma~\ref{cc:uf:trunc:lem:coefficients}(c) give
\begin{equation}\label{cc:uf:mean:thresholds}
 \ell_*^m\bigl(|K|_{m,1/2}+|\Omega|_{m,1/2}\bigr)
 \le C_m\mathsf D^C\ell_*^{1/2}\le C_m/A\qquad(m\ge1).
\end{equation}
Label derivatives preserve the endpoint and divergence constraints,
the time independence of the divergence-free field, and the class of
label gradients in \eqref{cc:uf:mean:integrated}. Theorem~\ref{cc:uf:mean:inverse}
bounds the map $L\mapsto v$ by $A$ in $C_t\Holder{0}$. The term $Kv$
enters directly, and $\Omega v$ enters through time integration,
an operator of norm at most $T\le1$.
Thus \eqref{cc:uf:mean:thresholds} gives the coefficient bounds for
Lemma~\ref{cc:uf:trunc:lem:lifting} with $C_{\rm low}=A$.
Applying its H\"older estimate to the source $L=-\int_0^tg(t')\dd t'$
of \eqref{cc:uf:mean:integrated} gives
\begin{equation}\label{cc:uf:mean:holder-bound}
 \norm v_{[n],1/2}\le C_n\mathsf D^C\norm g_{[n],1/2}.
\end{equation}
By \eqref{cc:uf:trunc:eq:holder-sup} and \eqref{cc:uf:hier:eq:coefficient-bound},
$\|D_aX\|_{[n],1/2}\le C_n\ell_*^{-1/2}\mathsf D^C\le C_n\mathsf D^{C'}$,
with $C'$ absolute. Using $B=D_aXv$ and $g=(D_aX)^T\bar f$ in
\eqref{cc:uf:mean:holder-bound}, the product estimate \eqref{cc:uf:trunc:eq:product}
gives $\|B\|_{[n],1/2}\le C_n\mathsf D^C\|\bar f\|_{[n],1/2}$.
The comparison \eqref{cc:uf:trunc:eq:holder-sup} now gives \eqref{cc:uf:mean:bound}.
\end{proof}
